\section{Chapter Summary}
\label{sec:ch2-summary}
This chapter established the conceptual and scholarly foundation for the thesis.
It defined the core concepts of adaptive reading: the closed sensing-to-intervention loop, micro-interventions and just-in-time adaptation, context preservation with recovery metrics such as Reading Resume Time, the condition of central vision loss that motivates the work, the oculomotor events (fixations, saccades, and regressions) extracted from gaze, and typography as the medium through which interventions act.
The review of related work positioned this thesis within the lineage of prior Reading the Reader projects \cite{refsgaard2024rtr, pereira2024typography, palinec2024reactive} and within the broader literature on eye-tracking in reading research and on intervention design.
The market survey then showed that no existing product combines real-time gaze sensing, automatic typographic adaptation, and researcher-facing experiment control in a single modular platform.

Taken together, the literature and the market landscape indicate that an adaptive reading system requires robust gaze-processing pipelines, carefully designed intervention mechanisms, measurable context-preservation outcomes, and a modular, experiment-ready architecture.
The gap analysis distilled these into three concrete opportunities, concerning experimentability, the sensing layer, and the integration of external decision modules, which the final problem statement in Sec.~\ref{sec:final-problem-statement} converts into the research questions that guide the remainder of the thesis.
\Cref{ch:methodology} describes the methodology adopted to address them.
